\chapter{Margin, Liquidation and Loss Allocation}\label{ch:m3:margin-liquidation-and-loss-allocation}

On 10 October 2025, after the announcement of a 100\% tariff on Chinese imports into the United States,
about USD~19 billion of leveraged crypto positions were liquidated within twenty-four hours, and open interest in
\omterm{def:m3:perpetual-futures:perp}{perpetual futures} across major venues fell by 43\%, from USD~217 billion to 123 billion. Bitcoin fell
to about USD~106{,}500, and the \omterm{def:m3:blockchains-for-traders:key}{tokens} CoinDesk Research tracked fell by about 47\% on average. On some venues the
liquidations were so large and so fast that the buffers meant to absorb their losses ran out, and the
venues closed profitable positions of other traders, who had done nothing wrong, to balance their books.
A crypto derivatives venue has no clearing members standing between it and its customers, and no default
fund contributed by banks: its customers' own margin, its \omterm{def:m3:margin-liquidation-and-loss-allocation:engine}{liquidation engine}, its \omterm{def:m3:margin-liquidation-and-loss-allocation:fund}{insurance fund} and, at
the end, its winning customers are the whole default waterfall. This chapter follows a leveraged position
from its margin to its \omterm{def:m3:margin-liquidation-and-loss-allocation:prices}{liquidation price}, through the engine that closes it, and down the waterfall; then
it builds and dissects a cascade.

\section{Isolated, cross and portfolio margin}

On a regulated futures exchange a customer's margin is set by the clearing house and collected through a
clearing member (One Quant Book 1, chapter 20). A crypto venue sets and collects margin itself, from each
account, in the account's own collateral, continuously.

\begin{definition}[Isolated margin, cross margin]\label{def:m3:margin-liquidation-and-loss-allocation:modes}
Under \emph{isolated margin}\index{isolated margin} each position has its own margin, assigned by the trader;
its losses can consume only that margin, and it is liquidated alone. Under \emph{cross
margin}\index{cross margin} all positions of an account share the account's collateral; gains on one support
losses on another, and the account is liquidated when its total equity falls below the sum of its
maintenance margins.
\end{definition}

\omterm{def:m3:margin-liquidation-and-loss-allocation:modes}{Isolated margin} caps what a single bet can lose and makes its \omterm{def:m3:margin-liquidation-and-loss-allocation:prices}{liquidation price} easy to compute; \omterm{def:m3:margin-liquidation-and-loss-allocation:modes}{cross margin}
uses collateral more efficiently and survives a loss on one position if others gain, at the risk of losing
everything in the account at once. Portfolio margining (One Quant Book 1, chapter 20), offered by some venues
to large accounts, sets the requirement from the scenario loss of the whole portfolio, including spot holdings
used as collateral. Maintenance margin is usually tiered: the rate rises with the size of the position, because
a larger position takes longer and costs more to liquidate.

\begin{dated}{2026-09}{Tiered maintenance margin}\label{dat:m3:margin-liquidation-and-loss-allocation:tiers}
Binance computes maintenance margin as the position's notional value times the maintenance margin rate of its
bracket, minus a maintenance amount that makes the requirement continuous from one bracket to the next;
brackets, rates and maximum leverage are published per contract and adjusted. BitMEX's instrument data on 24
September 2026 gave XBTUSD an initial margin of 1\% and a maintenance margin of 0.5\% at the base risk limit.
\end{dated}

\section{The liquidation price}

\begin{definition}[Liquidation price, bankruptcy price]\label{def:m3:margin-liquidation-and-loss-allocation:prices}
The \emph{liquidation price}\index{liquidation price} of a position is the \omterm{def:m3:perpetual-futures:index}{mark price} at which its margin plus
unrealised P\&L falls to its maintenance margin, so that the venue takes it over. Its \emph{bankruptcy
price}\index{bankruptcy price} is the price at which margin plus unrealised P\&L is exactly zero.
\end{definition}

\begin{proposition}[Liquidation prices of isolated positions]\label{prop:m3:margin-liquidation-and-loss-allocation:liq}
Let a linear position of $q$ units be entered at $E$ with margin $M$ and maintenance margin rate $m$ (no
maintenance amount). A long ($q > 0$) is liquidated at and goes bankrupt at
\[
P_{\mathrm{liq}} = \frac{qE - M}{q(1 - m)}, \qquad P_{\mathrm{bk}} = E - \frac{M}{q},
\]
and a short of $n = -q$ units at $P_{\mathrm{liq}} = (M + nE)/(n(1 + m))$ and $P_{\mathrm{bk}} = E + M/n$. With
leverage $\Lambda = qE/M$ the long's \omterm{def:m3:margin-liquidation-and-loss-allocation:prices}{liquidation price} is $E(1 - 1/\Lambda)/(1 - m)$. A long inverse position of
$N$ one-dollar contracts with margin $M$ in coin is liquidated at $N(1 + m)/(M + N/E)$.
\end{proposition}

\begin{proof}
For the long, $M + q(P - E) = m\,qP$ gives the first formula and $M + q(P - E) = 0$ the second; the short's are
the same with the sign of the P\&L reversed. For the inverse long, the equity in coin is $M + N(1/E - 1/P)$ and the
maintenance requirement $mN/P$; equating them gives the stated price.
\end{proof}

The symbol $\Lambda$ is local to this chapter (in One Quant Book 4 it denotes a compensator). At 10 times leverage a
long is liquidated about 9.5\% below its entry, at 50 times about 1.5\% below. The inverse long is liquidated sooner
than the linear one at the same leverage, because its coin margin loses value in the same move: it is long twice, as
\cref{ch:m3:perpetual-futures} explained. \Cref{fig:m3:margin-liquidation-and-loss-allocation:liqprice} draws the three.

\begin{omfigure}
\begin{tikzpicture}
\begin{axis}[width=11cm,height=5cm,xmode=log,log basis x=10,
  xlabel={leverage},ylabel={liquidation price, \% of entry},
  tick label style={font=\small},label style={font=\small},
  xmin=2,xmax=100,ymin=40,ymax=160,grid=major,grid style={gray!25},
  xtick={2,5,10,20,50,100},xticklabels={2,5,10,20,50,100},
  legend columns=3,legend style={font=\footnotesize,at={(0.5,-0.3)},anchor=north,draw=none,fill=none,/tikz/every even column/.append style={column sep=8pt}},
  table/col sep=comma]
\addplot[omThm,very thick] table[x=leverage,y=long_linear]{figdata/markets-3/18-margin-liquidation-and-loss-allocation/liqprice.csv};\addlegendentry{long, linear}
\addplot[omDef,very thick,dashed] table[x=leverage,y=short_linear]{figdata/markets-3/18-margin-liquidation-and-loss-allocation/liqprice.csv};\addlegendentry{short, linear}
\addplot[omExo,very thick,dotted] table[x=leverage,y=long_inverse]{figdata/markets-3/18-margin-liquidation-and-loss-allocation/liqprice.csv};\addlegendentry{long, inverse}
\end{axis}
\end{tikzpicture}

\omcaption{\omterm{def:m3:margin-liquidation-and-loss-allocation:prices}{Liquidation prices} of isolated positions as a function of leverage, with a maintenance margin rate of
0.5\% (\cref{prop:m3:margin-liquidation-and-loss-allocation:liq}). A long inverse position with coin margin is
liquidated sooner than a linear one: at 2 times leverage, after a 33\% fall instead of 50\%. Data: the chapter's
tutorial.}\label{fig:m3:margin-liquidation-and-loss-allocation:liqprice}
\end{omfigure}

Two consequences matter for risk. Liquidation is triggered by the \omterm{def:m3:perpetual-futures:index}{mark price} (\cref{ch:m3:perpetual-futures}), so a
trader's \omterm{def:m3:margin-liquidation-and-loss-allocation:prices}{liquidation price} is about the venue's index, not the trades he sees; and the gap between the liquidation
and \omterm{def:m3:margin-liquidation-and-loss-allocation:prices}{bankruptcy prices}, the maintenance margin, is all the venue has to close the position without loss.

\section{Liquidation engines}

\begin{definition}[Liquidation engine]\label{def:m3:margin-liquidation-and-loss-allocation:engine}
A \emph{liquidation engine}\index{liquidation engine} is the venue's system that monitors every account against
its maintenance margin at the \omterm{def:m3:perpetual-futures:index}{mark price}, takes over the positions of accounts that breach it, and closes them,
usually by sending orders into the venue's own order book, sometimes by transferring them to designated
liquidity providers.
\end{definition}

The engine is the venue's clearing house in software. It cancels the account's open orders, may first try a
partial liquidation that reduces the position to a lower margin tier, then sells (for a long) into the bids. If it
sells above the \omterm{def:m3:margin-liquidation-and-loss-allocation:prices}{bankruptcy price}, what is left of the trader's margin is kept, on many venues, by the \omterm{def:m3:margin-liquidation-and-loss-allocation:fund}{insurance
fund}; if below, the fund pays the difference. The engine is fast, mechanical and indifferent to price: its orders
are exactly the forced selling of the liquidity spirals of One Quant Book 1, chapter 31, arriving when the book is thinnest.
FTX's help pages described the alternative: backstop liquidity providers, invited from its top volume tiers, kept at least USD~500{,}000
on the venue and pledged to absorb at least USD~0.1 million of liquidations a minute and 0.3 million an hour, taking over an account's
position and collateral before it went bankrupt. And the engine's venue can itself fail when it is needed: on 13 March 2020 BitMEX was
hit by two denial-of-service attacks, at 02:16 and 12:56 UTC, which delayed or prevented requests to the platform.

\section{Insurance funds, auto-deleveraging and socialised losses}

\begin{definition}[Insurance fund]\label{def:m3:margin-liquidation-and-loss-allocation:fund}
An \emph{insurance fund}\index{insurance fund} is a pool of assets a derivatives venue holds to absorb the losses
of positions liquidated below their \omterm{def:m3:margin-liquidation-and-loss-allocation:prices}{bankruptcy prices}, funded mostly by the remaining margin of positions
liquidated above them.
\end{definition}

\begin{dated}{2026-09}{An insurance fund balance}\label{dat:m3:margin-liquidation-and-loss-allocation:fund}
BitMEX's public insurance endpoint reported, for 23 September 2026, a bitcoin \omterm{def:m3:margin-liquidation-and-loss-allocation:fund}{insurance fund} of about 3{,}700
bitcoin and a USDT fund of about 30.6 million USDT. Venues publish their funds' balances in different ways and at
different frequencies.
\end{dated}

\begin{definition}[Auto-deleveraging, socialised loss]\label{def:m3:margin-liquidation-and-loss-allocation:adl}
\emph{Auto-deleveraging}\index{auto-deleveraging} is the closing of profitable opposite positions of other traders,
at the bankrupt position's \omterm{def:m3:margin-liquidation-and-loss-allocation:prices}{bankruptcy price} or at the mark, when the \omterm{def:m3:margin-liquidation-and-loss-allocation:fund}{insurance fund} cannot absorb a liquidation's
loss, with those traders selected by a published ranking. A \emph{socialised loss}\index{socialised loss} is a
loss spread across a group of traders, for example all profitable accounts pro rata to their gains, rather than
imposed on those selected by a ranking.
\end{definition}

\begin{omfigure}
\centering
\begin{tikzpicture}[font=\small,
  w/.style={draw,thick,rounded corners=2pt,align=center,minimum height=0.85cm,text width=11.5cm},
  arr/.style={-{Stealth[length=2.2mm]},thick}]
\node[w,draw=omDef,fill=omDef!8] (a) at (0,4.4) {1.\ the trader's own margin, down to the bankruptcy price};
\node[w,draw=omThm,fill=omThm!8] (b) at (0,3.0) {2.\ the liquidation engine sells into the book: surplus or deficit};
\node[w,draw=omProp,fill=omProp!8] (c) at (0,1.6) {3.\ the insurance fund pays the deficit};
\node[w,draw=omExo,fill=omExo!8] (d) at (0,0.2) {4.\ auto-deleveraging: profitable opposite positions closed, by rank};
\node[w,draw=gray,fill=gray!8] (e) at (0,-1.2) {5.\ (on some venues, or historically) losses socialised across winners};
\draw[arr] (a) -- (b); \draw[arr] (b) -- (c); \draw[arr] (c) -- (d); \draw[arr] (d) -- (e);
\end{tikzpicture}

\omcaption{The default waterfall of a crypto derivatives venue. There are no clearing members and no default fund
contributed by them, as in One Quant Book 1, chapter 5: after the loser's margin and the venue's \omterm{def:m3:margin-liquidation-and-loss-allocation:fund}{insurance fund}, the
losses fall on the winners. Schematic.}\label{fig:m3:margin-liquidation-and-loss-allocation:waterfall}
\end{omfigure}

\begin{dated}{2026-09}{Auto-deleveraging rules at two venues}\label{dat:m3:margin-liquidation-and-loss-allocation:adlrules}
Binance applies \omterm{def:m3:margin-liquidation-and-loss-allocation:adl}{auto-deleveraging} only if its \omterm{def:m3:margin-liquidation-and-loss-allocation:fund}{insurance funds} are unable to accept a bankrupt position; profitable
opposite positions are ranked by P\&L percentage times effective leverage and closed at the \omterm{def:m3:margin-liquidation-and-loss-allocation:prices}{bankruptcy price} of the
liquidated order, without trading fees, and each position shows an indicator of its place in the queue. Hyperliquid's
documentation triggers \omterm{def:m3:margin-liquidation-and-loss-allocation:adl}{auto-deleveraging} when an account's or an isolated position's value becomes negative, ranks
profitable users by (\omterm{def:m3:perpetual-futures:index}{mark price} over entry price) times (notional over account value), and closes their positions at
the previous \omterm{def:m3:perpetual-futures:index}{mark price} against the underwater user.
\end{dated}

The ranking chooses who pays: the most profitable and most leveraged winners first, because their closure removes
the most exposure for the least number of accounts. To the winner, \omterm{def:m3:margin-liquidation-and-loss-allocation:adl}{auto-deleveraging} is an involuntary exit at a price
he did not choose, often in the middle of the move he was positioned for, and his hedges elsewhere are left open. A
market maker short the perpetual against spot can find its short closed and its spot long left naked in a crash. After
October 2025, an analysis by Tarun Chitra of Gauntlet, reported by Protos, argued that the queue-based ranking common
to many venues closed about USD~650 million more of profitable positions on Hyperliquid than the minimum needed, and
proposed a risk-aware pro-rata rule instead.

\section{A liquidation cascade dissected}

\begin{definition}[Liquidation cascade]\label{def:m3:margin-liquidation-and-loss-allocation:cascade}
A \emph{liquidation cascade}\index{liquidation cascade} is a sequence in which forced sales by \omterm{def:m3:margin-liquidation-and-loss-allocation:engine}{liquidation engines}
move the price enough to reach further positions' \omterm{def:m3:margin-liquidation-and-loss-allocation:prices}{liquidation prices}, whose forced sales move it further.
\end{definition}

Whether a price shock stops or cascades depends on how many units sit at each \omterm{def:m3:margin-liquidation-and-loss-allocation:prices}{liquidation price} and on how far each
unit sold moves the price.

\begin{proposition}[The local cascade multiplier]\label{prop:m3:margin-liquidation-and-loss-allocation:kappa}
Let each unit sold move the price down by $k$ (a linear, permanent impact), and let $\rho(P)$ be the number of units
whose \omterm{def:m3:margin-liquidation-and-loss-allocation:prices}{liquidation prices} lie in a unit interval around $P$. Then $\kappa(P) = k\,\rho(P)$ is the further fall caused,
through liquidations, by one point of fall near $P$. Where $\kappa$ is constant and below one, a fall of $\delta$ ends
as a fall of $\delta/(1 - \kappa)$; where $\kappa \ge 1$, the cascade does not stop until the price leaves the region.
The final price does not depend on the order or the batch size in which triggered positions are liquidated.
\end{proposition}

\begin{proof}
A fall $\delta$ triggers $\rho\delta$ units, which cause a fall $\kappa\delta$, which triggers $\kappa^2\delta$, and so
on: the geometric series sums to $\delta/(1 - \kappa)$ if $\kappa < 1$ and diverges otherwise. With linear permanent
impact, the price after any set of liquidations is the shocked price minus $k$ times the units sold, whatever their
order; liquidating every position whose price has been reached, in any batches, stops at the same largest fixed point.
\end{proof}

\begin{omfigure}
\begin{tikzpicture}
\begin{axis}[width=11cm,height=5cm,ybar,bar width=4pt,
  xlabel={price (entry prices around 100)},ylabel={$\kappa$},
  tick label style={font=\small},label style={font=\small},
  xmin=59,xmax=100,ymin=0,ymax=1.3,grid=major,grid style={gray!25},
  table/col sep=comma]
\addplot[fill=omThm!60,draw=omThm] table[x=price,y=kappa]{figdata/markets-3/18-margin-liquidation-and-loss-allocation/kappa.csv};
\draw[omDef,very thick,dashed] (axis cs:59,1) -- (axis cs:100,1);
\end{axis}
\end{tikzpicture}

\omcaption{The local cascade multiplier of \cref{prop:m3:margin-liquidation-and-loss-allocation:kappa} for the
tutorial's 4{,}000 leveraged longs (USD~167 million of notional, leverage from 2 to 50 times) and an impact of 0.12\%
of the price per USD~1 million sold; the dashed line marks $\kappa = 1$. Between about 92 and 96 the multiplier exceeds one: a cascade that reaches that
zone runs through it. Synthetic positions. Data: the chapter's tutorial.}\label{fig:m3:margin-liquidation-and-loss-allocation:kappa}
\end{omfigure}

\Cref{fig:m3:margin-liquidation-and-loss-allocation:drawdown} shows the consequence. Small shocks are amplified
two- to threefold and stop. A shock of 1.11\% is enough to carry the price into the zone where $\kappa > 1$, and the
cascade then runs to below 90: a ninefold amplification. Capping leverage at ten times removes the cluster of \omterm{def:m3:margin-liquidation-and-loss-allocation:prices}{liquidation
prices} near the entry and raises the threshold to 5.37\%; doubling the depth of the book halves $\kappa$, which then stays below one
everywhere, so the fall is amplified smoothly, about twofold, and reaches 10\% only after a shock of 5.55\%; without
forced selling there is no cascade at all. The tutorial checks that the result does not depend on the size of the
engine's batches, as the proposition says, and ablates each mechanism in turn.

\begin{omfigure}
\begin{tikzpicture}
\begin{axis}[width=11cm,height=5cm,
  xlabel={initial shock, \%},ylabel={total fall, \%},
  tick label style={font=\small},label style={font=\small},
  xmin=0,xmax=10,ymin=0,ymax=26,grid=major,grid style={gray!25},
  legend columns=2,legend style={font=\footnotesize,at={(0.5,-0.3)},anchor=north,draw=none,fill=none,/tikz/every even column/.append style={column sep=8pt}},
  table/col sep=comma]
\addplot[omThm,very thick] table[x=shock,y=base]{figdata/markets-3/18-margin-liquidation-and-loss-allocation/drawdown.csv};\addlegendentry{base case}
\addplot[omDef,very thick,dashed] table[x=shock,y=cap10]{figdata/markets-3/18-margin-liquidation-and-loss-allocation/drawdown.csv};\addlegendentry{leverage capped at 10x}
\addplot[omExo,very thick,dotted] table[x=shock,y=depth2]{figdata/markets-3/18-margin-liquidation-and-loss-allocation/drawdown.csv};\addlegendentry{book twice as deep}
\addplot[gray,thick] table[x=shock,y=noforced]{figdata/markets-3/18-margin-liquidation-and-loss-allocation/drawdown.csv};\addlegendentry{no forced selling}
\end{axis}
\end{tikzpicture}

\omcaption{Total fall after the cascade as a function of the initial shock, for the base case and three ablations. In
the base case the fall jumps from under 4\% to over 10\% between shocks of 1.0\% and 1.25\%; capping leverage at
ten times moves the jump past 5\%, and a book twice as deep removes it. Synthetic positions and book. Data: the chapter's
tutorial.}\label{fig:m3:margin-liquidation-and-loss-allocation:drawdown}
\end{omfigure}

A gap is worse than a slide. When the price gaps down 5\% at once, as on a surprise announcement, every position whose
\omterm{def:m3:margin-liquidation-and-loss-allocation:prices}{liquidation price} lies inside the gap is closed at the gapped price, below many \omterm{def:m3:margin-liquidation-and-loss-allocation:prices}{bankruptcy prices}: in the tutorial,
USD~3.21 million of deficits that the \omterm{def:m3:margin-liquidation-and-loss-allocation:fund}{insurance fund} must pay, and the price ends near 84. The weekend problem follows the
loss down the waterfall.

\section{Tutorial: liquidation prices and a cascade}\label{tut:m3:margin-liquidation-and-loss-allocation:cascade}

\begin{tutorial}
\textbf{Goal.} Derive \omterm{def:m3:margin-liquidation-and-loss-allocation:prices}{liquidation prices}, simulate a cascade of leveraged longs selling into a book of stated depth, test
its stability under a finer time step, and ablate each mechanism. \textbf{End state:}
\cref{fig:m3:margin-liquidation-and-loss-allocation:liqprice,fig:m3:margin-liquidation-and-loss-allocation:kappa,fig:m3:margin-liquidation-and-loss-allocation:drawdown}
and the numbers of the text.
\begin{tutsteps}
\item \textbf{Liquidation and \omterm{def:m3:margin-liquidation-and-loss-allocation:prices}{bankruptcy prices}.} \Cref{prop:m3:margin-liquidation-and-loss-allocation:liq} in code.
\omcode{code/firm/liquidation/firm_liquidation.py}{26}{37}{Liquidation and bankruptcy prices of an isolated linear position.}\label{lst:m3:margin-liquidation-and-loss-allocation:liq}
\item \textbf{The cascade.} Triggered positions are sold into the book, all at once or in batches.
\omcode{code/markets-3/18-margin-liquidation-and-loss-allocation/python/m3_liq.py}{59}{83}{A liquidation cascade with a configurable batch size.}\label{lst:m3:margin-liquidation-and-loss-allocation:cascade}
\item \textbf{Stability and ablations.} Run \texttt{cascade} with \texttt{batch} of 100, 10 and 1 and compare with the
batch of everything triggered; then set \texttt{forced\_selling=False}, cap leverage and halve the impact.
\item \textbf{Run} \texttt{fig\_liq.py} for the chart data.
\end{tutsteps}
\textbf{What to change next.} Replace the linear impact by a square-root impact and check whether the batch size still
does not matter; let the book refill between batches; add short positions and a funding payment.
\end{tutorial}

\section{Build: the margin and liquidation engine}\label{bld:m3:margin-liquidation-and-loss-allocation:liquidation}

\begin{build}
\bfield{Purpose} The miniature firm must know, for every position on every venue, how far it is from liquidation and
what a liquidation would cost; and its risk simulations need the venue's engine and waterfall reproduced.
\bfield{Interface} \texttt{Tier(max\_notional, mmr, maint\_amount)}, \texttt{maintenance\_margin(notional, tiers)};
\texttt{liq\_price\_linear}, \texttt{bankruptcy\_price\_linear}, \texttt{liq\_price\_inverse};
\texttt{cross\_health(wallet, positions, tiers)}; \texttt{sell\_into\_book(qty, bids)}; \texttt{adl\_rank(accounts,
price)}; \texttt{liquidate(account, bids, fund, opposite, price)}.
\bfield{Rules} \omterm{def:m3:perpetual-futures:index}{Mark prices} from \texttt{firm.perp}; tiers continuous at their boundaries; a liquidation's surplus goes to
the fund and its deficit is paid by it; if the fund cannot pay, the position is closed against opposite positions in ADL
order at the \omterm{def:m3:margin-liquidation-and-loss-allocation:prices}{bankruptcy price}, leaving book and fund unchanged.
\bfield{Acceptance tests} \texttt{code/firm/liquidation/tests/}: continuity of tiers; linear and inverse \omterm{def:m3:margin-liquidation-and-loss-allocation:prices}{liquidation
prices} against the closed forms; cross-margin health; selling into a book; surplus to the fund; ADL in rank order.
\bfield{Stretch} Partial liquidation to a lower tier; backstop liquidity providers; a pro-rata ADL rule; cross-venue
exposure to one underlying.
\end{build}

\begin{omsources}
\item CoinDesk Research, ``Market Spotlight: Inside Crypto's \$19 Billion Liquidation Event'', 17 October 2025.
\item Binance, ``What Is Auto-Deleveraging (ADL)'' and ``Leverage and Margin of USDS-M Futures'', support pages;
Hyperliquid documentation, ``Auto-deleveraging''. Accessed September 2026.
\item BitMEX, public API insurance and instrument endpoints, 23--24 September 2026; BitMEX blog, ``How We Are Responding to the 13 March
DDoS Attacks'', March 2020.
\item FTX, help centre, ``Backstop Liquidity Provider Program'' (updated 14 September 2022), Internet Archive copy.
\item Protos, ``Outdated algorithm caused \$650M excess losses on Hyperliquid, report'', 10 December 2025.
\end{omsources}

\section{Exercises}

\begin{exercise}[$\star$]\label{exo:m3:margin-liquidation-and-loss-allocation:1}
A 20-times long at 100 with a maintenance margin rate of 0.5\%: what are its liquidation and \omterm{def:m3:margin-liquidation-and-loss-allocation:prices}{bankruptcy prices}?
\end{exercise}

\begin{exercise}[$\star$]\label{exo:m3:margin-liquidation-and-loss-allocation:2}
A trader holds one position at 50 times leverage in \omterm{def:m3:margin-liquidation-and-loss-allocation:modes}{cross margin} with a large unused balance. Is she liquidated 2\% below
entry? Explain.
\end{exercise}

\begin{exercise}[$\star$]\label{exo:m3:margin-liquidation-and-loss-allocation:3}
Name the layers of a crypto venue's default waterfall, in order.
\end{exercise}

\begin{exercise}[$\star\star$]\label{exo:m3:margin-liquidation-and-loss-allocation:4}
Compare the \omterm{def:m3:margin-liquidation-and-loss-allocation:prices}{liquidation prices} of a linear and an inverse long at 5 times leverage and 0.5\% maintenance margin. Why do
they differ?
\end{exercise}

\begin{exercise}[$\star\star$]\label{exo:m3:margin-liquidation-and-loss-allocation:5}
Near the current price $\kappa = 0.8$. A news shock moves the price down 1\%. How far does it fall in all, if $\kappa$
stays constant?
\end{exercise}

\begin{exercise}[$\star\star$]\label{exo:m3:margin-liquidation-and-loss-allocation:6}
A market maker is long spot and short the perpetual. How can \omterm{def:m3:margin-liquidation-and-loss-allocation:adl}{auto-deleveraging} hurt it, and what would it do about it?
\end{exercise}

\begin{exercise}[$\star\star\star$]\label{exo:m3:margin-liquidation-and-loss-allocation:7}
\emph{Coding.} With \texttt{critical\_shock}, find the threshold for leverage caps of 20 and 10 times. What does this
suggest about venue limits on leverage?
\end{exercise}

\begin{exercise}[$\star\star\star$]\label{exo:m3:margin-liquidation-and-loss-allocation:8}
\emph{Find the flaw.} ``A cascade simulation that liquidates positions one at a time is more accurate than one that
liquidates everything triggered at once, so its answer will be different.''
\end{exercise}

\section{Problem: Anatomy of a Cascade}

\begin{problem}[{Weekend problem --- from a shock to auto-deleveraging}]\label{pb:m3:margin-liquidation-and-loss-allocation:1}
Use the tutorial's population of 4{,}000 leveraged longs (USD~166.6 million of notional), a maintenance margin rate of 0.5\%,
an impact of 0.0012 price points per unit sold (0.12\% per USD~1 million at a price of 100) and an \omterm{def:m3:margin-liquidation-and-loss-allocation:fund}{insurance fund} of USD~2
million.

\medskip\noindent\textbf{Part I --- One position.}
\begin{enumerate}
\item A long of 250 units at 100 with 25 times leverage: what are its margin, \omterm{def:m3:margin-liquidation-and-loss-allocation:prices}{liquidation price} and \omterm{def:m3:margin-liquidation-and-loss-allocation:prices}{bankruptcy price}?
\item If it is liquidated at 95.5, what does the \omterm{def:m3:margin-liquidation-and-loss-allocation:fund}{insurance fund} pay?
\item And if at 96.8?
\item Why does the fund receive money in the second case?
\item Why does the venue use the \omterm{def:m3:perpetual-futures:index}{mark price}, not the last trade, to trigger it?
\end{enumerate}

\medskip\noindent\textbf{Part II --- The cascade.}
\begin{enumerate}[resume]
\item What is the smallest initial shock after which the price falls by 10\% or more in all?
\item What does $\kappa$ exceed between about 92 and 96, and what does that mean?
\item What are the thresholds with leverage capped at 10 times, and with a book twice as deep?
\item Why is the final price the same whatever the batch size?
\item What would break that invariance?
\end{enumerate}

\medskip\noindent\textbf{Part III --- The waterfall.}
\begin{enumerate}[resume]
\item After a 5\% gap, where does the price end, and what are the deficits?
\item How much reaches \omterm{def:m3:margin-liquidation-and-loss-allocation:adl}{auto-deleveraging}?
\item Who is deleveraged first under a ranking by P\&L percentage times leverage?
\item What does a deleveraged market maker lose, and what does it keep?
\item How would a pro-rata rule change who pays?
\end{enumerate}

\medskip\noindent\textbf{Part IV --- Judgement.}
\begin{enumerate}[resume]
\item Should venues cap leverage?
\item Is an \omterm{def:m3:margin-liquidation-and-loss-allocation:fund}{insurance fund} of USD~2 million enough for this market?
\item How should a trading firm use a venue's ADL indicator?
\item State the \emph{named result}: the critical shock, and the loss that reaches the \omterm{def:m3:margin-liquidation-and-loss-allocation:fund}{insurance fund} and then
\omterm{def:m3:margin-liquidation-and-loss-allocation:adl}{auto-deleveraging} after a 5\% gap.
\item In one sentence: who stands behind a crypto derivatives trade?
\end{enumerate}
\end{problem}

\section{Interview questions}

\begin{interviewq}[$\star$ \iqroles{trader}]\label{iq:m3:margin-liquidation-and-loss-allocation:1}
Derive the \omterm{def:m3:margin-liquidation-and-loss-allocation:prices}{liquidation price} of a 10-times long.
\end{interviewq}

\begin{interviewq}[$\star$ \iqroles{risk}]\label{iq:m3:margin-liquidation-and-loss-allocation:2}
What is \omterm{def:m3:margin-liquidation-and-loss-allocation:adl}{auto-deleveraging}, and why should a firm that is not leveraged care about it?
\end{interviewq}

\begin{interviewq}[$\star\star$ \iqroles{researcher}]\label{iq:m3:margin-liquidation-and-loss-allocation:3}
What determines whether a price shock turns into a \omterm{def:m3:margin-liquidation-and-loss-allocation:cascade}{liquidation cascade}?
\end{interviewq}

\begin{interviewq}[$\star\star$ \iqroles{developer}]\label{iq:m3:margin-liquidation-and-loss-allocation:4}
How would you design a \omterm{def:m3:margin-liquidation-and-loss-allocation:engine}{liquidation engine} to minimise losses to the \omterm{def:m3:margin-liquidation-and-loss-allocation:fund}{insurance fund}?
\end{interviewq}

\begin{interviewq}[$\star\star$ \iqroles{risk, trader}]\label{iq:m3:margin-liquidation-and-loss-allocation:5}
Compare isolated and \omterm{def:m3:margin-liquidation-and-loss-allocation:modes}{cross margin} for a market maker quoting 50 perpetuals.
\end{interviewq}

\begin{interviewq}[$\star\star\star$ \iqroles{researcher, risk}]\label{iq:m3:margin-liquidation-and-loss-allocation:6}
Using public data, how would you estimate where the \omterm{def:m3:margin-liquidation-and-loss-allocation:prices}{liquidation prices} of a venue's open interest lie?
\end{interviewq}
